% ============================================================
%  第 2 章：总体分析（问题分析、数据说明、建模流程）
% ============================================================
\section{总体分析}

\subsection{问题分析}

（对各子问题进行定性分析，指出每个问题的本质类型（如预测问题、优化问题、评价问题、分类问题），并初步确定可能的建模方向。）

\textbf{问题一分析：}（该问题本质上是一个...问题，需要解决...，可选用的方法有...）

\textbf{问题二分析：}（分析思路同上）

\textbf{问题三分析：}（分析思路同上）

\subsection{数据说明}

（如题目提供了数据，在此说明数据的来源、格式、规模、是否存在缺失值/异常值等基本情况。若题目未提供数据，此节可省略。）

数据来源：（说明数据来自题目附录文件 xxx.xlsx / xxx.csv）

数据规模：共 $n$ 条记录，$m$ 个变量/特征，其中...

数据质量：（说明缺失值情况、异常值情况）

\subsection{整体建模流程}

本文的整体建模流程如图 \ref{fig:flowchart} 所示。

\begin{figure}[H]
    \centering
    % \includegraphics[width=0.8\textwidth]{flowchart.png}
    \fbox{\parbox{0.6\textwidth}{\centering\vspace{3cm}（在此插入流程图）\vspace{3cm}}}
    \caption{整体建模流程图}
    \label{fig:flowchart}
\end{figure}

首先进行数据预处理（数据清洗、标准化等），然后依次针对三个子问题建立模型并求解，最后对模型进行检验和评价。
